\documentclass[11pt]{article}

\usepackage[margin=1in]{geometry}
\usepackage{booktabs}
\usepackage{graphicx}
\usepackage{hyperref}
\usepackage{amsmath}

\title{Respiratory Disease-Group Classification from Auscultation Audio}
\author{Draft}
\date{\today}

\begin{document}

\maketitle

\begin{abstract}
This manuscript studies respiratory disease-group prediction from auscultation
audio. The target labels are Normal, Airway disease, and Lung Parenchymal
disease. The experimental focus compares direct three-class classification
against a two-stage cascade that first separates Normal from Abnormal and then
separates Airway from Lung Parenchymal disease. The runnable implementation
currently includes Audio Spectrogram Transformer (AST), Whisper-based, and
ResNet50-based backbones under a shared evaluation contract.
\end{abstract}

\section{Introduction}

Respiratory auscultation provides clinically meaningful acoustic evidence for
screening and grouping pulmonary conditions. This study targets three disease
groups: Normal, Airway disease, and Lung Parenchymal disease. The primary
research question is whether a direct three-class model or a two-stage cascade
offers a better operational structure for disease-group prediction.

\section{Dataset}

The current disease-group experiments use wav files organized by label
directory. The primary final labels are \texttt{Normal}, \texttt{Airway}, and
\texttt{Lung\_Parenchymal}. The split-local \texttt{Abnormal} directory is used
as a binary overlay for cascade stage 1 and represents the union of
\texttt{Airway} and \texttt{Lung\_Parenchymal}; it is not a fourth final label.

Future revisions should report patient-level and clip-level counts, fold
membership, test-set counts, annotation coverage, and any exclusion criteria
after the dataset catalog has been regenerated and reviewed.

\section{Methods}

\subsection{Direct Three-Class Classification}

The direct approach trains one backbone-specific classifier with a softmax
output over \texttt{Normal}, \texttt{Airway}, and
\texttt{Lung\_Parenchymal}. Every input clip passes through the same model and
receives one final label.

\subsection{Cascade Classification}

The cascade approach trains two binary classifiers per backbone. Stage 1 predicts
\texttt{Normal} versus \texttt{Abnormal}. Clips predicted as
\texttt{Abnormal} are routed to stage 2, which predicts \texttt{Airway} versus
\texttt{Lung\_Parenchymal}. Clips predicted as \texttt{Normal} by stage 1 do
not reach stage 2.

\subsection{Backbone-Specific Acoustic Frontends}

Because this study evaluates transfer-learning pipelines, each pretrained
backbone is used with the acoustic frontend expected by its pretraining
configuration rather than forcing all models to share one spectrogram
representation. AST therefore uses an AST-style Kaldi fbank frontend, Whisper
uses a Whisper-style log-mel frontend, and ResNet50 uses an image-style
spectrogram frontend. This design preserves the input distribution assumed by
each pretrained encoder. Comparability is enforced at the disease-group task
level through shared label definitions, fold splits, checkpoint selection
rules, evaluation metrics, and direct-versus-cascade reporting, not by treating
the frontend as interchangeable across backbones. Consequently, the model
comparison should be interpreted as a comparison of pretrained backbone
pipelines with native acoustic frontends rather than as a pure encoder
architecture ablation.

\subsection{AST Pipeline}

The current runnable implementation uses clip-level AST classification. Audio
is resampled to 16 kHz, converted to AST-style fbank features, and fed to a
Hugging Face \texttt{ASTModel} encoder. The current AST disease-group configs
use 128 mel bins and a pretrained-compatible maximum feature length of 1024
frames. The classifier head is implemented in the repository and supports
linear or MLP heads. Binary tasks use BCE or focal loss, while multi-class
tasks use cross entropy.

\subsection{Whisper Pipeline}

The Whisper-based implementation uses a local Whisper-like audio encoder and
OpenAI Whisper encoder checkpoints. Audio is resampled to 16 kHz and converted
to 80-bin log-mel features over a 30 s frame context. The default Whisper
encoder context is 1500 after the stride-2 audio convolution. Binary tasks use
one positive-class logit so that threshold optimization, checkpoint evaluation,
and cascade routing share the same contract as the AST pipeline. The encoder
can be kept frozen or partially fine-tuned through the experiment config. In
the partial L1 setting, only the final Whisper encoder block is trainable after
pretrained checkpoint loading. Multi-class tasks use class logits with cross
entropy.

\subsection{ResNet50 Pipeline}

The ResNet50 implementation uses torchvision ResNet50 pretrained on ImageNet
by default. Audio is resampled to 16 kHz and trimmed or padded to 15.0 s. The
frontend computes 128-bin log-mel power maps with 25 ms windows and 10 ms hop
length, applies HPSS on the mel power representation, and stacks the full,
harmonic, and percussive log-mel maps as three image channels. The stacked
representation is resized to 224 by 224 pixels and normalized with ImageNet
mean and standard deviation. The ResNet50 fully connected layer is replaced by
an identity layer, and the classifier is attached directly to the pooled
2048-dimensional feature. The encoder can be frozen, partially fine-tuned, or
fully fine-tuned. Partial fine-tuning unfreezes the final residual stages in a
fixed order, while the initial convolution and batch normalization layers are
trainable only in the full fine-tuning setting.

\section{Experimental Design}

The initial experiment set contains three cross-validation tasks per runnable
backbone:

\begin{itemize}
  \item Direct: \texttt{Normal\_vs\_Airway\_vs\_LungParenchymal}
  \item Cascade stage 1: \texttt{Normal\_vs\_Abnormal}
  \item Cascade stage 2: \texttt{Airway\_vs\_LungParenchymal}
\end{itemize}

For Whisper and ResNet50, the first comparison arm uses frozen encoders as the
transfer-learning baseline. A second partial L1 arm uses the same tasks,
splits, labels, and frontend policy, but unfreezes the final encoder block for
Whisper or the final residual stage for ResNet50. These runs are stored under
separate result roots so that frozen and partially fine-tuned checkpoints are
not mixed during evaluation.

For the final held-out comparison, the same test clips should be evaluated by
the direct model and the cascade. The cascade comparison should report both
final-label performance and stage-specific routing outcomes.

\section{Evaluation Metrics}

Evaluation should report accuracy, precision, recall, specificity, balanced
accuracy, F1 score, ROC AUC, PR AUC, brier score, and confusion matrix. For the
three-class final disease-group task, macro-averaged metrics and per-class
support should be emphasized. For cascade analysis, routing error categories
should be reported alongside final-label confusion matrices.

\section{Model Implementation Status}

AST, Whisper, and ResNet50 are implemented as runnable disease-group backbones.
They share the same dataset contract, label mapping, checkpoint metadata,
metrics, and reporting outputs, while preserving backbone-specific frontends
appropriate for transfer learning.

\section{Discussion}

The direct model offers a single decision boundary over the final label space,
while the cascade model encodes a clinically interpretable Normal-versus-
Abnormal gate. The cascade can reduce the burden on the fine-grained disease
classifier, but routing errors at stage 1 can block disease clips from stage 2
or force true Normal clips into a disease subtype. These structural tradeoffs
should guide both quantitative evaluation and qualitative error analysis.

\section{Reproducibility Notes}

The repository currently stores configs for AST, Whisper, and ResNet50 direct
three-class CV, cascade stage 1 CV, and cascade stage 2 CV. The AST direct
direct and cascade configs write to an AST family result root while preserving
the canonical task identifiers. Whisper and ResNet50 configs write under
family-specific frozen-baseline result roots. Additional Whisper and ResNet50
partial L1 configs write under separate partial L1 roots. Result directories
are intentionally ignored by git, so final paper numbers should be regenerated
from documented commands and archived separately when the experiments are
finalized.
Lightweight CV summary scripts may be used to audit evaluated folds before
transferring final numbers into the manuscript. The preferred cascade config
files are named by stage and task so the Normal-versus-Abnormal gate and
Airway-versus-Lung Parenchymal subtype classifier are not confused during
experiment execution.
The AST direct three-class five-fold CV run is currently treated as an active
experiment; numerical results should be inserted only after all folds complete
and checkpoint evaluation outputs have been audited.
The ResNet50 direct-to-cascade five-fold CV queue was started on
June 20, 2026. ResNet50 numerical results should likewise be inserted only
after all folds complete and checkpoint evaluation outputs have been audited.

\end{document}
